% ===========================================================================
\section{Related Work}
\label{sec:related_work}
% ===========================================================================

\subsection{National inventory methods}

National agricultural inventories generally follow the 2006 IPCC Guidelines
\citep{ipcc2006} and, increasingly, the 2019 Refinement \citep{ipcc2019}, which updated
several agricultural emission factors and introduced climate-specific disaggregation for
soil nitrous oxide.

The IPCC framework defines several methodological tiers. Tier~1 combines default emission
factors with national activity data. Tier~2 replaces selected defaults with
country-specific parameters or more detailed activity characterisation. Tier~3 uses
higher-order models or direct national measurement systems.

Pakistan's national agricultural inventories have been prepared mainly with Tier~1 methods
and national activity totals \citep{gcisc,mir2022}. This approach supports comprehensive
national reporting but provides limited information about subnational patterns or
individual processes.

\subsection{Global gridded inventories}

Several international datasets provide independent estimates of agricultural emissions for
Pakistan. FAOSTAT calculates national agricultural and land-use emissions from FAO activity
data using mainly IPCC Tier~1 methods and reports results by category and gas
\citep{faostat,tubiello2013}.

EDGAR provides national totals and gridded emissions at approximately \SI{0.1}{\degree}
resolution using a globally consistent methodology \citep{edgar2025booklet,crippa2020};
this paper's comparison table keeps one row from the prior 2024-vintage booklet
\citep{edgar2024} because the 2025 release revised Pakistan's agricultural
\CHfour{} history and the revision itself is informative. Climate
TRACE reports country- and subsector-level agricultural emissions using a combination of
inventories, models and observation-informed methods \citep{climatetrace}.

These products provide useful external reference points, but they are not substitutes for a
nationally compiled inventory. They generally use globally harmonised or model-derived
assumptions and may not incorporate the most appropriate national activity data or system
definitions. In this study, they are used as external benchmarks rather than as reported
national values.

\subsection{Spatial allocation of emissions}

Allocating national or subnational emissions to a grid is a dasymetric-mapping problem. An
emission total is distributed according to a spatial proxy that represents the location or
intensity of the underlying activity. Similar approaches are widely used in fossil-CO$_2$
inventories and global emissions gridmaps \citep{oda2018,crippa2020}.

Several open datasets provide relevant agricultural proxies. The Gridded Livestock of the
World (\GLW{}) dataset provides modelled species-specific livestock densities
\citep{glw4,gilbert2018}. The Spatial Production Allocation Model (MapSPAM) provides
crop-specific physical areas \citep{mapspam2020}. ESA WorldCover provides a
\SI{10}{\metre} land-cover classification \citep{worldcover2021}, while NASA FIRMS and
other satellite products identify active fires and burned areas \citep{firms,giglio2016}.

The quality of a spatial inventory depends on whether each proxy represents the process
being allocated. Livestock density is appropriate for emissions produced directly by
animals, paddy area for rice methane, cropland for fertiliser-related emissions, and
satellite fire occurrence for the location of residue burning. A single proxy applied to an
entire inventory category may misallocate emissions when that category contains several
distinct activity streams.

A further requirement is mass conservation. Some spatial inventories distribute emissions
without subsequently rescaling the mapped values to the original control total. In such
cases, the gridded product may no longer reproduce the national or provincial inventory.
The framework developed here treats reconciliation as an invariant: every mapped emission
stream must sum to its corresponding control total.

\subsection{Agricultural ammonia inventories}

Agricultural ammonia inventories have largely developed within the air-quality community,
separately from greenhouse-gas inventories used for climate reporting. The EMEP/EEA
Guidebook represents livestock ammonia through a stage-resolved mass flow. Total ammoniacal
nitrogen is tracked through grazing, housing, yards, storage and field application, with
emission factors varying by animal category and manure-management system
\citep{emep2023manure}.

This process resolution is needed because ammonia losses occur at specific stages of manure
handling. By contrast, Tier~1 greenhouse-gas methods often represent manure nitrogen using
annual fractions that aggregate across several stages. Combining the two frameworks
therefore requires careful alignment of their system boundaries and nitrogen pools.

\citet{sanh2026} provide a regional reference for South Asia. Their framework estimates
agricultural \NHthree{} emissions for the SAARC countries across five sectors using an
explicit nitrogen-flow model and provides an openly available \SI{0.1}{\degree} grid
\citep{sanh2026data}. We use this inventory to compare both the national magnitude and the
spatial pattern of our ammonia estimates.

Other global products, including EDGAR and HTAP, also cover Pakistan. However, they are
built without national activity data and rely partly on the same underlying FAO statistics
used by many national and regional inventories. They therefore provide useful bounds and
methodological comparisons, but agreement with them does not constitute fully independent
confirmation.

\subsection{Crop-residue burning inventories}

Crop-residue burning can be estimated using two broad approaches. Bottom-up inventories
calculate the mass burned from crop production, residue-to-product ratios, dry-matter
fractions, combustion factors, and the proportion of residues burned. National
greenhouse-gas inventories generally use this approach, and the burned fraction is its
dominant uncertainty.

Top-down or observation-based approaches use satellite fire detections, fire radiative
power, or burned-area products. These datasets provide much stronger information about the
location and timing of fires, but conversion to burned mass depends on additional
assumptions and also carries substantial uncertainty.

This study uses the two types of information for separate purposes. The national quantity
burned is calculated using a bottom-up inventory, while satellite fire detections are used
only to allocate that quantity spatially. Satellite detections therefore determine the
mapped distribution but do not provide an independent estimate of the national total, a
distinction that matters because the two are sometimes conflated.

\subsection{Contribution relative to previous work}

Compared with conventional national inventory workbooks, this framework integrates
calculation, uncertainty analysis, spatial allocation, external evaluation and historical
reconstruction in a reproducible codebase. Compared with global gridded products, it uses
nationally grounded activity data and province-level control totals within a hybrid
Tier~1/Tier~2 agricultural inventory.

Its main contribution is the combination of these elements. The national inventory
establishes the emission totals, process-specific proxies determine their spatial
distribution, and exact reconciliation keeps the provincial, gridded and district products
consistent with the reported figures. To our knowledge, this combination has not previously
been implemented for Pakistan's agricultural sector.
